% ECONOMICS_PROBLEM: {"id": "D82-650", "title": "Markov persuasion regret--violation tradeoffs: resolved gap", "jel_code": "D82", "source_id": "OP-0181"}
% !TeX program = xelatex
\documentclass[11pt,a4paper]{article}
\usepackage[margin=23mm]{geometry}
\usepackage{fontspec}
\setmainfont{Latin Modern Roman}
\usepackage{amsmath,amssymb,mathtools}
\usepackage{xcolor,enumitem,booktabs,longtable,array,xurl,hyperref}
\definecolor{muted}{HTML}{53636C}
\hypersetup{colorlinks=true,urlcolor=blue,linkcolor=blue}
\setlength{\parindent}{0pt}
\setlength{\parskip}{5pt}
\newcommand{\R}{\mathbb R}
\newcommand{\E}{\mathbb E}
\newcommand{\N}{\mathbb N}
\newcommand{\ind}{\mathbf 1}
\newcommand{\Var}{\operatorname{Var}}
\newcommand{\Cov}{\operatorname{Cov}}
\newcommand{\supp}{\operatorname{supp}}
\DeclareMathOperator*{\argmin}{arg\,min}
\DeclareMathOperator*{\argmax}{arg\,max}
\newcommand{\entrymeta}[3]{{\sffamily\small\color{muted}\textbf{\texttt{#1}}\quad|\quad #2\par #3\par}}
\newcommand{\Problemref}[1]{\texttt{#1}}
\newcommand{\JELref}[2]{\texttt{#2} (JEL \texttt{#1})}
\begin{document}
\section*{Markov persuasion regret--violation tradeoffs: resolved gap}
\textbf{Problem ID:} \texttt{D82-650}\quad \textbf{JEL:} \texttt{D82}\par
% BEGIN ECONOMICS STATEMENT
\label{problem:OP-0181}
% GAME_THEORY_PROBLEM: {"local_id": "GT-051", "original_number": 51, "kind": "H", "entry_kind": "statement_only", "published": false, "review_required": true, "category": "game_theory"}
\label{game:GT-051}
{\sffamily\small\color{muted}
\textbf{GT-051} | Learning and information design | \textbf{H}: the cited paper closes the advertised exponent gap.
\par}
\subsection*{Definitions and mathematical statement}
An episodic Markov persuasion process has finite layer sets $X_0,\ldots,X_L$, finite outcomes $\Omega$ and actions $A$, unknown outcome kernels $\mu(\omega\mid x)$, unknown transition kernel $P(x'\mid x,\omega,a)$, and unknown mean rewards $r_S,r_R\in[0,1]$. Bounded reward vectors are drawn independently each episode with these means. The initial state and finite sets are known. In episode $t$ a learner chooses a recommendation rule $\phi_t(a\mid x,\omega)$; recommendations are followed during learning. The learner observes reached states, outcomes, and the two realized rewards for the chosen action, without seeing rewards for other actions. Write $q_t(x,\omega,a)$ for the resulting occupancy probability. For a reached recommendation, define
\[
 b_t(a,x)\in\operatorname*{arg\,max}_{b\in A}
 \sum_\omega\mu(\omega\mid x)\phi_t(a\mid x,\omega)r_R(x,\omega,b).
\]
For zero-probability recommendations its choice is immaterial. Let
\[
 U_S(\phi_t)=\sum_{x,\omega,a}q_t(x,\omega,a)r_S(x,\omega,a),
\]
\[
 V_T=\sum_{t=1}^T\sum_{x,\omega,a}q_t(x,\omega,a)
 [r_R(x,\omega,b_t(a,x))-r_R(x,\omega,a)].
\]
The sum over outcomes for each recommendation is nonnegative by best-response optimality. A persuasive rule has zero such violation at every reached recommendation. If $U_S^*$ is optimal value over persuasive rules under the true process, regret is $R_T=T U_S^*-\sum_tU_S(\phi_t)$. The original catalogue asks for matching bounds outside $\alpha\in[1/2,2/3]$. In the source's partial-feedback model, the advertised tradeoff is already established throughout
\[
 \alpha\in[1/2,1],\qquad
 R_T=\widetilde O(T^\alpha),\qquad V_T=\widetilde O(T^{1-\alpha/2}),
\]
with the corresponding simultaneous little-$o$ improvement ruled out by Theorem 5. Here $\widetilde O$ suppresses logarithmic factors and factors depending on the fixed finite process dimensions.
\subsection*{Short English statement}
The alleged lower-bound exponent gap is not open in the cited Markov persuasion model.
\subsection*{Variants and scope boundaries}
The source assumes recommendation obedience during learning and records violations relative to myopic best responses; it does not assume every exploratory policy is persuasive. Different feedback or strategically learning receivers define separate problems.
\subsection*{Source relationship and status}
Section 5.3 compares its full-interval result with a narrower guarantee in older work. The catalogue mistakenly assigns that older gap to the cited paper. This entry is preserved as a historical correction, not replaced by an invented new open question.
\subsection*{References}
\begingroup\footnotesize\raggedright
\textbf{[S1]} Francesco Bacchiocchi, Francesco Emanuele Stradi, Matteo Castiglioni, Alberto Marchesi, and Nicola Gatti (2024). \emph{Markov Persuasion Processes: Learning to Persuade from Scratch}. arXiv:2402.03077v2. \textit{Verified locator:} Sections 2--3; Section 5.2, Theorems 3--4; Section 5.3, Theorem 5 and comparison paragraph. \url{https://arxiv.org/html/2402.03077v2}
\par\endgroup

\subsection*{Common conventions: Source mathematical conventions}

\textbf{Finite games.} For a finite set $A$, $\Delta(A)$ is its probability
simplex. Players choose independent mixed strategies in Nash equilibrium;
correlation is allowed only when explicitly part of the solution concept.
In a game with payoff functions $u_i$ and strategy profile $x$, an additive
$\varepsilon$ Nash equilibrium satisfies
\[
 u_i(x)\geq u_i(x_i',x_{-i})-\varepsilon
 \quad\text{for every player $i$ and every admissible deviation $x_i'$}.
\]
For numerical approximation, payoffs are normalized as locally specified.
An $\varepsilon$ payoff residual is distinct from distance at most
$\varepsilon$ to the set of exact equilibria. Point convergence, convergence
of empirical play, and convergence in distance to an equilibrium set are
also distinct.

\textbf{Computational representations.} Unless stated otherwise, a complexity
question takes rational data with binary encoding; polynomial time is measured
in total bit length. A fixed-accuracy polynomial algorithm is not necessarily
polynomial in $1/\varepsilon$, and neither is necessarily polynomial in
$\log(1/\varepsilon)$. Succinct games and value, demand, or sampling oracles
must declare their interfaces. Exact real solutions require an explicit
representation; an unspecified list of arbitrary real numbers is not a
finite computational output. A claimed algorithmic barrier is meaningful
only for its specified model and reductions.

\textbf{Dynamic games.} Histories, information, observation, and allowable
randomization are part of the model. A stationary strategy depends only on
the current state; a positional strategy in a deterministic graph game
chooses one action at each controlled vertex. Nash equilibrium at an initial
state and equilibrium after every history are different requirements.
An expectation of a pathwise $\liminf$ payoff, a $\liminf$ of expected averages,
a limiting expected average, and a uniform finite-horizon equilibrium are
different criteria. None is silently substituted for another.

\textbf{Allocation and mechanisms.} A complete allocation assigns every item
exactly once unless otherwise stated. Zero-valued goods, negative values,
capacity constraints, feasibility, lotteries, and copies of goods affect
fairness definitions and are specified locally. Dominant-strategy incentive
compatibility, Bayesian incentive compatibility, universal truthfulness,
and truthfulness in expectation impose different inequalities. Whenever
an objective ratio has zero denominator, its convention is stated or those
instances are excluded.

\textbf{Infinite-dimensional models.} Expectations, integrals, stochastic
equations, and optimization objectives require their local regularity,
integrability, filtrations, and admissible domains. A backward--forward
mean-field system is a boundary-value problem; it is not an initial-value
semiflow without an additional construction. Long-time, large-population,
and vanishing-noise limits require an order of limits and a convergence
topology. If a minimum need not be attained, the objective is its infimum
or supremum and the approximation requirement is stated separately.
% END ECONOMICS STATEMENT
\end{document}
